\chapter{Material e Métodos}
\section {Considerações Iniciais}
Essa é a área "Material e Métodos".

Nela, você deve discriminar o que foi adotado até o momento, ou prováveis métodos a serem usados para solucionar o problema. Indicar como e onde os dados serão obtidos


Abaixo, um exemplo de tabela, possivelmente útil nessa etapa do trabalho \ref{tb:tabela_dados}.

\begin{table}[htbp]
\caption{Exemplo de dados em tabelas}
\label{tb:tabela_dados}
\centering
\setlength{\tabcolsep}{5pt}
\begin{tabular}{cccccc}
\hline
Coluna 1  &Coluna 2  &Coluna 3 \\
L2Col1 &L2Col2 &L2Col3 \\
\hline
Dado1C1 &100 &200 \\
Dado2C1 &300 &400 \\
\hline
\textbf{Total} &\textbf{400} &\textbf{600} \\
\hline
\end{tabular}
\\
\text{\footnotesize Fonte: O autor}
\end{table}

\section{Exemplo de Código Fonte}
\subsection{ Teste de subseção}
Esse é um exemplo de código fonte ~\ref{code:exemplo_fonte}, que também pode ser útil nessa etapa do projeto de pesquisa.

%\lstset{frame=Trbl,numbers=left}
\begin{lstlisting}[caption={Exemplo de código fonte}, language=Python, label=code:exemplo_fonte]
	# Python program to check if the input number
	# is odd or even.
	# A number is even if division by 2 gives a
	# remainder of 0.
	# If the remainder is 1, it is an odd number.
	
	num = int(input("Enter a number: "))
	if (num % 2) == 0:
		print("{0} is Even".format(num))
	else:
		print("{0} is Odd".format(num))
\end{lstlisting}	

